%!TEX root = ../main.tex

% Main Text
\section{Introduction}

\noindent Autonomous field robots are reshaping precision agriculture, replacing manual scouting and hand-steered passes with continuous, real-time operation between crop rows. The foundational competence on which most of these platforms rest is also the most deceptively simple: follow the rows. A forward-facing camera observes the canopy, a perception module localizes the rows, and a \emph{guidance line} is constructed that summarizes where the rows go; a steering controller then tracks this line to keep the vehicle centred between rows~\cite{cth1,cth2}. Because the guidance line is the single interface between perception and actuation, its accuracy --- expressed not as a detection score but as a heading in degrees and a cross-track offset in centimetres --- directly determines whether the robot stays in the lane, damages plants, or requires human intervention.

\noindent A dominant and edge-friendly recipe for this task pairs a lightweight object detector with a downstream line fit: the detector localizes row segments as bounding boxes at successive distances, and a line is fitted through the box centroids from the bottom of the frame toward the vanishing point~\cite{diao_yolox}. This family is attractive for agricultural robots because the detector can be exported to integer precision and run on commodity edge accelerators. Yet it inherits a perspective-geometry pathology that is specific to navigation. Near-field rows occupy large, texture-rich image regions and are localized precisely; far-field rows shrink to a few pixels, crowd together, and blur into soil and neighbouring vegetation, so their box edges are estimated with large and \emph{spatially varying} uncertainty. Critically for steering, it is exactly these unreliable far-field detections that anchor the upper end of the guidance line and therefore set the look-ahead heading. A single mislocalized far box can rotate the fitted line and inject a heading error that the controller integrates forward in time~\cite{cth2}.

\noindent We argue that the persistent difficulty here is not primarily a detection-architecture problem but a \emph{measurement and estimation} problem, and that two gaps explain why prior work --- across detection-and-line-fit~\cite{diao_yolox}, segmentation-and-scan~\cite{cth1}, row-anchor regression~\cite{alnet}, and curve regression~\cite{rowdetr} pipelines --- leaves accuracy on the table. First, the optimization target and the deployed outcome are disconnected: detectors are trained and selected on mean average precision (mAP), a proxy that is never causally tied to the quantity the robot actually steers on, the guidance-line heading. Improving mAP does not guarantee a better line, and a worse mAP does not imply worse navigation, because mAP averages over all boxes equally while navigation is dominated by a few far-field ones. Second, regardless of the perception paradigm, the guidance line is fitted with \emph{equal} (or heuristic) weights: every point contributes identically even though the far points are an order of magnitude noisier than the near points. The estimator thus spends its precision where it is least reliable. Box-regression uncertainty has been modelled before to improve \emph{detection} --- via a Gaussian edge head~\cite{gaussianyolo}, a KL localization loss~\cite{klloss}, or the distributional regression head used in modern YOLO detectors~\cite{gfl} --- but it has not been read out as a \emph{calibrated} per-detection covariance and propagated into the guidance-line fit, where it would matter most for steering.

\noindent We close both gaps with \textbf{CALM-Row} (Calibrated Aleatoric Localization for Maize-row guidance), a post-head module that turns a stock YOLO detector into a navigation-native, uncertainty-aware crop-row guidance estimator without changing the detector graph. The key observation is that the information needed to down-weight unreliable far boxes is already present, for free, in the detector. The Distribution Focal Loss (DFL) regression head~\cite{gfl} predicts, per anchor, a probability distribution over $R$ bins for each box edge $e\in\{l,t,r,b\}$; the box coordinate is the distribution mean $\mu_e=\sum_k k\,p_e(k)$, while its \emph{variance} $\sigma_e^2=\sum_k (k-\mu_e)^2 p_e(k)$ --- exploited by GFLv2~\cite{gflv2} only as a scalar quality score, and otherwise discarded by the standard box decode --- is a direct, learned, per-edge localization-uncertainty signal. Sharp distributions (confident near boxes) give small variance; flat distributions (ambiguous far boxes) give large variance. CALM-Row reads this distribution, converts it to a per-box centroid covariance by linear error propagation, and uses the resulting precision $w_i=1/\sigma_{c_x,i}^2$ to weight a differentiable generalized-least-squares (GLS) line fit. Our contributions are as follows. \textbf{(i)} We extract a per-detection localization-uncertainty signal \emph{for free} from the existing DFL distribution --- no extra head, no extra learnable parameters in the training losses --- and propagate it through closed-form error propagation to a per-box centroid covariance ($\sigma_{c_x}^2=(s^2/4)(\sigma_l^2+\sigma_r^2)$, with $s$ the stride). \textbf{(ii)} We fold this covariance into a differentiable, uncertainty-weighted GLS guidance-line fit ($x=ay+b$, parameterized along the image rows to avoid the vertical-line degeneracy of near-vertical crops), whose closed-form normal equations and parameter covariance also yield a principled heading $\theta=\arctan a$ and an associated heading uncertainty. \textbf{(iii)} We train the detector with two auxiliary, parameter-free losses that operate purely on the DFL distribution: a distance-aware calibration loss (DDC), a Gaussian negative-log-likelihood that ties the predicted centroid variance to the realised squared centroid error on matched anchors so the uncertainty becomes \emph{calibrated} rather than merely relative, and a line-stability loss (LSL) whose gradients flow through the weights $w_i$ into the DFL logits, teaching the network to down-weight unreliable far boxes for a stabler look-ahead line. \textbf{(iv)} We design an evaluation in navigation units --- heading error (degrees) as the controller-independent primary metric, line-fit RMSE (px and cm via a stated homography), and a static geometric cross-track proxy at the look-ahead row (a closed-loop pure-pursuit simulation over video sequences is left to future work) --- alongside calibration (ECE, reliability) and edge latency, with generalization measured in-domain on CRDLD and as cross-dataset transfer to CRBD, and statistical claims supported by multiple seeds, paired significance tests, and effect sizes. An illustrative synthetic mechanism check is given in Section~\ref{sec:res-synth}; all dataset-level navigation, calibration, mAP and latency numbers are reported only after the experiments of Section~\ref{sec:results} are run.

\noindent Because all of CALM-Row's computation is post-head CPU arithmetic over post-NMS boxes, it is detector- and neck-agnostic and adds \emph{zero} new operators to the on-device graph: the exported INT8 model --- whether targeting an RKNN (Rockchip) NPU or TensorRT (Jetson) --- remains a vanilla YOLO, and the calibration and line-stability losses introduce no new learnable parameters during training. CALM-Row therefore upgrades the accuracy of the guidance line, and connects it to the navigation outcome, at no inference-time cost on the edge hardware that agricultural robots actually deploy.